\chapter{Fats}
\index{Fats}

\section*{Definition and Overview}
Fats, also called lipids, are an essential part of a healthy diet. They provide energy, help the body absorb fat-soluble vitamins, build cell membranes, and supply essential fatty acids that the body cannot make itself. However, the type and amount of fat eaten matter greatly for health: some fats protect the heart, while others, eaten in excess, raise cholesterol and increase the risk of heart disease and stroke. Understanding the different kinds of fat, the foods that contain them, and how much to eat is a cornerstone of preventing chronic disease. This chapter explains dietary fats, their roles, and practical guidance for eating them wisely.

\section*{Epidemiology}
Dietary fat intake varies widely between populations and individuals. In many countries, fats contribute roughly 30--35\% of total energy intake, close to the recommended range of 20--35\%. However, a large share of this comes from unhealthy saturated and trans fats, and high intakes of saturated fat are linked to elevated LDL (bad) cholesterol, which affects around one-third of adults. Heart disease remains a leading cause of death, and dietary change, including replacing saturated fats with unsaturated fats, is recognised as one of the most effective population-level strategies to reduce cardiovascular risk. Consumption of trans fats has fallen substantially in many countries following reformulation of processed foods, but saturated fat intake remains above recommended levels for many people.

\section*{Aetiology and Pathophysiology}
Different fats have different chemical structures and very different effects on the body.
\begin{itemize}
    \item \textbf{Saturated fats:} found mainly in animal products such as fatty meat, butter, cream, and full-fat dairy, and in coconut and palm oils; they raise LDL cholesterol and increase cardiovascular risk
    \item \textbf{Trans fats:} formed when liquid oils are hydrogenated to make them solid; found in some processed, fried, and baked foods; they raise LDL cholesterol, lower HDL (good) cholesterol, and are the most harmful type
    \item \textbf{Monounsaturated fats:} found in olive oil, avocado, nuts, and canola oil; they lower LDL cholesterol and are heart-protective
    \item \textbf{Polyunsaturated fats:} found in sunflower, safflower, and fish oils; including the omega-3 fatty acids in oily fish, which reduce inflammation and support heart and brain health
    \item \textbf{Cholesterol:} a waxy substance carried in the blood that is needed for cell membranes and hormones, but excess LDL in the blood contributes to plaque in arteries
    \item \textbf{Essential fatty acids:} omega-3 and omega-6 fats cannot be made by the body and must come from food
    \item \textbf{Energy density:} fat provides 37 kilojoules per gram, more than double that of carbohydrate or protein, so high-fat diets can drive weight gain
\end{itemize}

\section*{Clinical Features}
Dietary fats do not cause symptoms in themselves; their effects are long-term and measured in the blood and body.
\begin{itemize}
    \item Elevated LDL cholesterol and triglycerides on blood tests, often with no symptoms at all
    \item Excess body weight and obesity from overconsumption of energy-dense fats
    \item Accumulation of visceral fat around the organs, associated with metabolic disease
    \item Fatty liver, detectable on imaging, in people with high-fat, high-energy diets
    \item Long-term consequences: atherosclerosis, heart attacks, strokes, and peripheral artery disease
    \item Deficiency states are rare but possible with extremely low-fat diets, causing dry skin, poor growth in children, and fat-soluble vitamin deficiency
\end{itemize}

\section*{Differential Diagnosis}
When blood lipid results are abnormal, the interpretation considers more than diet alone.
\begin{itemize}
    \item Familial hypercholesterolaemia and other inherited lipid disorders, which raise cholesterol regardless of diet
    \item Hypothyroidism, diabetes, kidney disease, and liver disease, which affect lipid levels
    \item Medications, including some diuretics and steroids, which can raise lipids
    \item Excess alcohol, which raises triglycerides
    \item Pregnancy, which naturally increases cholesterol levels
    \item Weight and physical activity patterns, which influence lipid levels independently of fat intake
\end{itemize}

\section*{When to Seek Medical Care}
Most people do not need to see a doctor about dietary fat, but anyone with a family history of early heart disease, known high cholesterol, diabetes, or who is concerned about their cardiovascular risk should have a heart health check.

\textbf{Contact your local emergency services for:}
\begin{itemize}
    \item Chest pain, pressure, or tightness, especially with sweating or breathlessness
    \item Sudden weakness, numbness, facial droop, or difficulty speaking
    \item Sudden severe abdominal pain, which may indicate acute pancreatitis from very high triglycerides
    \item Collapse or loss of consciousness
\end{itemize}

\section*{Diagnosis and Investigations}
\begin{itemize}
    \item \textbf{Lipid profile:} a fasting blood test measuring total cholesterol, LDL, HDL, and triglycerides
    \item \textbf{Cardiovascular risk assessment:} combining lipid results with age, blood pressure, smoking, diabetes, and family history to estimate overall risk
    \item \textbf{Repeat testing:} lipid levels are variable, so abnormal results are usually confirmed before starting treatment
    \item \textbf{Thyroid, kidney, liver, and glucose tests:} to exclude secondary causes of abnormal lipids
    \item \textbf{Genetic testing:} for familial hypercholesterolaemia when cholesterol is very high or there is a strong family history
\end{itemize}

\section*{Management}
\begin{itemize}
    \item \textbf{Replace saturated with unsaturated fats:} use olive, canola, and nut oils instead of butter and lard, and choose nuts, seeds, avocado, and oily fish
    \item \textbf{Choose lean proteins:} skinless chicken, fish, legumes, and reduced-fat dairy
    \item \textbf{Limit processed foods:} reduce biscuits, pastries, chips, fried fast food, and other sources of saturated and trans fats
    \item \textbf{Eat oily fish twice a week:} salmon, sardines, and mackerel provide heart-protective omega-3 fats
    \item \textbf{Watch portion sizes:} even healthy fats are energy-dense, so moderate portions of nuts, oils, and avocados are important
    \item \textbf{Maintain healthy weight and activity:} combined with diet, this is the most effective way to improve lipid levels
    \item \textbf{Medication:} statins and other lipid-lowering medicines are added when cardiovascular risk is high or lipids remain elevated despite lifestyle change
\end{itemize}

\section*{Complications}
\begin{itemize}
    \item Heart attack, stroke, and peripheral artery disease from long-term raised LDL cholesterol
    \item Acute pancreatitis from very high triglyceride levels
    \item Obesity, type 2 diabetes, and fatty liver disease from excess energy intake
    \item Gallstones, which are more common with high-fat diets and obesity
    \item Complications of statin therapy are uncommon but include muscle aches and, rarely, liver effects, which are monitored
\end{itemize}

\section*{Prognosis and Outcomes}
Dietary changes that replace saturated and trans fats with unsaturated fats reliably lower LDL cholesterol, and the effect on cardiovascular risk is well established: even modest reductions in blood cholesterol translate into substantially fewer heart attacks and strokes across a population. Weight loss and physical activity add further benefit. For people with familial hypercholesterolaemia, who have very high cholesterol from birth, medication is essential regardless of diet. With a balanced approach that includes healthy fats rather than avoiding fat altogether, most people can achieve durable improvements in lipids and long-term heart health.

\section*{Prevention and Self-Care}
\begin{itemize}
    \item Follow the local Dietary Guidelines, which recommend eating mostly unsaturated fats and limiting saturated fat
    \item Cook with small amounts of olive or canola oil, and grill, steam, or bake instead of frying
    \item Read food labels and compare saturated fat content per 100 g
    \item Include nuts, seeds, avocado, and oily fish in weekly meals
    \item Limit takeaway, processed snacks, and foods with hydrogenated or partially hydrogenated oils
    \item Ask a clinician how often you need cardiovascular risk assessment; timing depends on age, risk factors, and local guidance
    \item Stay active and maintain a healthy weight alongside dietary change
\end{itemize}

\section*{Sources and Further Reading}
The following reputable sources provide further information for healthcare students and patients seeking reliable, current advice about dietary fats.
\begin{itemize}
    \item National Health and Medical Research Council. \textit{Australian Dietary Guidelines.} https://www.eatforhealth.gov.au/
    \item Heart Foundation. \textit{Fats, oils, and heart health.} https://www.heartfoundation.org.au/
    \item Healthdirect Australia. \textit{Cholesterol and fats.} https://www.healthdirect.gov.au/
    \item Better Health Channel. \textit{Fats and oils in food.} https://www.betterhealth.vic.gov.au/
    \item NHS. \textit{Fat: the facts.} https://www.nhs.uk/live-well/eat-well/
    \item World Health Organization. \textit{Healthy diet and fats guidance.} https://www.who.int/
\end{itemize}
